\subsection{Field-affine verification, timing, and residual checking}\label{sec:zkrs:affine}

\begin{definition}[Complete ideal affine verification]\label{def:zkrs:affine-checks}
For a fixed public instance and public core, the ideal checker
receives hidden bound values
$(\widehat w_j)_j,p,p_{\rm in},c,t,\rho,\sigma$
with the following tests.
\begin{enumerate}[label=\textup{(\roman*)}]
\item \textbf{Claim:}
$y'=y+p_{\rm in}$.
\item \textbf{Incoming:}
$\sum_vq_vm_v+y'=\sum_vq_vp_v+p_{\rm in}$.
\item \textbf{Outgoing:}
$\sum_ja_j(\theta)\widehat w_j
=B_m(\theta)+D_c(\theta)$.
\item \textbf{Trace:} $\Tr(U_k)=0$ for every $k$;
these use only public values.
\item \textbf{First anchored equation:}
$A+\rho+\sum_kW_k(z)t_k=0$.
\item \textbf{Second anchored equation:}
\[
 A^2+\sigma+\sum_kW_k(z)^2
       (U_k+V_k(p,c,\theta)+t_k)=0.
\]
\end{enumerate}
A separate underlying commitment relation must link the
hidden packed columns to the fixed Boolean witness table.
\end{definition}

\begin{lemma}[The claimed-value and incoming tests are affine]\label{lem:zkrs:claim-in-affine}
Once the public parameters and $(m,y')$ are fixed,
the claim and incoming equations are affine over $L$
in the committed secret variables
$p,p_{\rm in}$.
\end{lemma}
\begin{proof}
Rewrite the claim equation as
$p_{\rm in}=y'+y$.
Rewrite the incoming equation as
$\sum_vq_vp_v+p_{\rm in}=\sum_vq_vm_v+y'$.
In both expressions the coefficients and right sides
are public, and each hidden element occurs with
degree one, proving the assertion.
\end{proof}

\begin{lemma}[The outgoing test is affine]\label{lem:zkrs:out-affine}
Once the public $(m,\theta)$ and the high-axis
weights are fixed, the outgoing test is affine
over $L$ in the committed packed columns
$\widehat w$ and coefficients $c$.
\end{lemma}
\begin{proof}
Expand $D_c(\theta)=\sum_u\lambda_u(\theta)c_u$.
Then the outgoing relation reads
\[
 \sum_ja_j(\theta)\widehat w_j+
 \sum_u\lambda_u(\theta)c_u=B_m(\theta).
\]
Every coefficient and the right side are publicly
determined. The hidden values occur linearly.
\end{proof}

\begin{lemma}[Both anchored equations are affine in hidden variables]\label{lem:zkrs:anchor-affine}
For fixed public $(m,y',\theta,U,z,A)$,
the two anchored equations of
Definition~\ref{def:zkrs:affine-checks}
are affine over $L$ in $(p,c,t,\rho,\sigma)$.
No product of two committed secret elements appears.
\end{lemma}
\begin{proof}
The first anchored equation is
$A+\rho+\sum_kW_k(z)t_k=0$:
its coefficients $W_k(z)$ and $A$ are public.
For the second, $A^2$ and every $W_k(z)^2$
are public scalars. By
Lemma~\ref{lem:pack:V-affine},
$V_k(p,c,\theta)$ is $L$-linear in the
secret pad and correction variables.
Thus
$A^2+\sigma+\sum_kW_k(z)^2(U_k+V_k+t_k)=0$
is affine in $(p,c,t,\sigma)$.
In particular $\sigma=\rho^2$
does \emph{not} appear among the
affine verifier equations.
\end{proof}

\begin{definition}[Algebraic residuals at the second challenge]\label{def:zkrs:errors}
For arbitrary fixed hidden variables and public $(\theta,U)$,
define
\[
 \delta:=\sigma+\rho^2,\qquad
 e_k:=U_k+V_k(p,c,\theta)+t_k^2+t_k.
\]
These values are \emph{defined mathematically}
for soundness analysis; the verifier need not
compute them directly or ask for nonlinear proofs.
\end{definition}

\begin{lemma}[Accepted affine equations imply a residual root]\label{lem:zkrs:root-condition}
If both anchored equations pass, then
\[
 \boxed{\delta+\sum_ke_kW_k(z)^2=0.}
\]
\end{lemma}
\begin{proof}
From the first equation,
$A=\rho+\sum_kW_k(z)t_k$.
Squaring in characteristic two gives
$A^2=\rho^2+\sum_kW_k(z)^2t_k^2$.
Substitute this into the second affine equation
and collect the scalar term
$\rho^2+\sigma=\delta$
and the summands
$U_k+V_k+t_k^2+t_k=e_k$.
The remaining equality is the displayed equation.
\end{proof}

\begin{theorem}[Any nonzero residual is detected with probability at least $1-n/q$]\label{thm:zkrs:residual-detection}
Fix $\theta,p,c,t,\rho,\sigma,U$ before
an independent uniform second challenge $z\in L$.
If $(\delta,e)\neq(0,0)$, then
even for an arbitrary public $A$ chosen after seeing $z$,
\[
 \Pr_z[\text{both anchored affine equations pass}]
 \leq\min\{1,n/q\}.
\]
\end{theorem}
\begin{proof}
The fixed-before-$z$ hypothesis makes
$(\delta,e)$ independent of the challenge.
By Lemma~\ref{lem:zkrs:root-condition},
acceptance implies
$\delta+\sum_ke_kW_k(z)^2=0$.
The anchored residual root bound
Theorem~\ref{thm:binary:anchor-root}
bounds this event by $n/q$ for any
nonzero residual tuple. The probability
is also at most one.
\end{proof}

\begin{corollary}[No secret squaring check is needed for this algebraic soundness]\label{cor:zkrs:no-quadratic-secret-check}
In the ideal algebraic setting, an adversarial choice
$\sigma\ne\rho^2$ fixed before the second challenge
is accepted by both anchored checks with probability
at most $\min\{1,n/q\}$.
\end{corollary}
\begin{proof}
The inequality $\sigma\ne\rho^2$
is equivalent to $\delta\ne0$,
so $(\delta,e)\ne(0,0)$ regardless
of the vector $e$. Apply
Theorem~\ref{thm:zkrs:residual-detection}.
\end{proof}

\begin{lemma}[Public trace checks impose no secret multiplication]\label{lem:zkrs:public-traces}
The trace constraints $\Tr(U_k)=0$ can be
checked from public field elements alone.
\end{lemma}
\begin{proof}
Each $U_k$ is explicitly published.
The trace is the finite field expression
$\sum_{i=0}^{n-1}U_k^{2^i}$ by
Definition~\ref{def:binary:trace}.
Its computation accesses only the
public value $U_k$ and public field arithmetic.
\end{proof}

\begin{remark}[What zero hidden-quadratic constraints means]
Only the \emph{six displayed algebraic checks}
above are affine in their hidden field values;
the trace check is public and the second
equation involves public $A^2$.
The bit-packing/commitment binding,
an implementation's internal arithmetic,
and the ability to prove these affine
relations in zero knowledge are separate
assumptions. Consequently this statement
is not a count of all gates or all transmitted
field elements in a concrete end-to-end PCS.
\end{remark}
